\section{Metric Differentiation, Levi--Civita Connection, and Hessian Calculus}
\label{app:rigid-calculus}
\subsection{Projector and square-root differentials}
\begin{definition}
For $X=U\Sigma V^\top\in\mathcal K_{a,b}$,
\begin{align}
S_L(X)&=(XX^\top)^{1/2}=U\Sigma U^\top,&
S_R(X)&=(X^\top X)^{1/2}=V\Sigma V^\top,\\
K_X(Y)&=S_L(X)Y+YS_R(X),&
A_X&=K_X^\dagger,\quad A_X|_{T_X}=K_X|_{T_X}^{-1}.
\end{align}
The square roots are restricted to the constant-rank stratum.
\end{definition}

\begin{lemma}[Differential of the tangent projector]
For $Z\in T_X\mathcal M_r$, let
\begin{equation}
Z=[U,U_\perp]\begin{pmatrix}Z_{11}&Z_{12}\\Z_{21}&0\end{pmatrix}
[V,V_\perp]^\top .
\end{equation}
Then
\begin{align}
DP_U[Z]&=U_\perp Z_{21}\Sigma^{-1}U^\top
+U\Sigma^{-1}Z_{21}^\top U_\perp^\top,\\
DP_V[Z]&=V_\perp Z_{12}^\top\Sigma^{-1}V^\top
+V\Sigma^{-1}Z_{12}V_\perp^\top,\\
DP_X[Z](Y)
&=DP_U[Z]Y(I-P_V)+(I-P_U)YDP_V[Z],\\
\nn{DP_X[Z](Y)}
&\leq \frac4a\nn Z\nn Y .
\label{eq:rigid-projector-derivative}
\end{align}
\end{lemma}
\begin{proof}
For a differentiable rank-$r$ path $X(t)$,
\begin{align}
P_UX&=X,&
(I-P_U)\dot P_U X&=(I-P_U)\dot X,\\
P_U\dot P_UP_U&=0,&
(I-P_U)\dot P_U(I-P_U)&=0,\\
(I-P_U)\dot P_UU&=U_\perp Z_{21}\Sigma^{-1},\\
\dot P_U&=(I-P_U)\dot P_UP_U+
P_U\dot P_U(I-P_U).
\end{align}
Apply the same identities to $X^\top$. Differentiate
$P_XY=P_UY+YP_V-P_UYP_V$ and use
\begin{align}
\on{DP_U[Z]}&\leq2a^{-1}\nn{Z_{21}},&
\on{DP_V[Z]}&\leq2a^{-1}\nn{Z_{12}},\\
\nn{Z_{21}}+\nn{Z_{12}}&\leq\sqrt2\nn Z .
\end{align}
\end{proof}

\begin{theorem}[Exact square-root and metric differentials]
\label{thm:metric-derivative}
Write $D_L=DS_L(X)[Z]$, $D_R=DS_R(X)[Z]$. In the $U,V$ block bases,
\begin{align}
D_L&\equiv\begin{pmatrix}T_L&Z_{21}^\top\\Z_{21}&0\end{pmatrix},&
\Sigma T_L+T_L\Sigma&=Z_{11}\Sigma+\Sigma Z_{11}^\top,\\
D_R&\equiv\begin{pmatrix}T_R&Z_{12}\\Z_{12}^\top&0\end{pmatrix},&
\Sigma T_R+T_R\Sigma&=Z_{11}^\top\Sigma+\Sigma Z_{11},\\
(T_L)_{ij}
&=\frac{Z_{11,ij}\sigma_j+\sigma_i Z_{11,ji}}{\sigma_i+\sigma_j},\\
(T_R)_{ij}
&=\frac{Z_{11,ji}\sigma_j+\sigma_i Z_{11,ij}}{\sigma_i+\sigma_j},\\
DK_X[Z](Y)&=D_LY+YD_R,\\
\nn{DK_X[Z](Y)}&\leq \frac{4b}{a}\nn Z\nn Y .
\label{eq:rigid-DK}
\end{align}
\end{theorem}
\begin{proof}
Differentiate $S_L^2=XX^\top$, $S_R^2=X^\top X$:
\begin{align}
S_LD_L+D_LS_L&=ZX^\top+XZ^\top,\\
S_RD_R+D_RS_R&=Z^\top X+X^\top Z,\\
ZX^\top+XZ^\top
&\equiv\begin{pmatrix}
Z_{11}\Sigma+\Sigma Z_{11}^\top&\Sigma Z_{21}^\top\\
Z_{21}\Sigma&0
\end{pmatrix}.
\end{align}
The constant-rank square-root path has zero normal--normal first derivative.
Consequently
\begin{align}
\nn{T_L}&\leq\frac ba\nn{Z_{11}},&
\nn{D_L}^2&=\nn{T_L}^2+2\nn{Z_{21}}^2
\leq\frac{4b^2}{a^2}\nn Z^2,\\
\nn{D_R}&\leq\frac{2b}{a}\nn Z,&
\nn{D_LY+YD_R}&\leq(\nn{D_L}+\nn{D_R})\nn Y .
\end{align}
\end{proof}

\begin{proposition}[Pseudoinverse derivative on the moving tangent bundle]
For $P=P_X$, $K=K_X$, $A=K_X^\dagger$, and $\dot K=DK_X[Z]$,
\begin{align}
DA_X[Z]
&=-A\dot KA+A^2\dot K(I-P)+(I-P)\dot KA^2,\\
P\,DA_X[Z]\,P&=-A(P\dot KP)A .
\label{eq:rigid-pseudoinverse-derivative}
\end{align}
For the embedded connection $\nabla^E_ZY=P_XDY[Z]$,
\begin{align}
(\nabla^E_ZK)(Y)&=P_XDK_X[Z](Y),\\
(\nabla^E_ZA)(Y)&=-A_X(\nabla^E_ZK)(A_XY),\\
\nn{(\nabla^E_ZA)(Y)}
&\leq \frac{4b}{a^3}\nn Z\nn Y .
\label{eq:rigid-bundle-derivative}
\end{align}
\end{proposition}
\begin{proof}
In the orthogonal decomposition $T_X\oplus T_X^\perp$,
\begin{align}
K&=\begin{pmatrix}K_T&0\\0&0\end{pmatrix},&
A&=\begin{pmatrix}K_T^{-1}&0\\0&0\end{pmatrix},\\
\dot A_{TT}&=-K_T^{-1}\dot K_{TT}K_T^{-1},&
\dot A_{TN}&=K_T^{-2}\dot K_{TN},&
\dot A_{NT}&=\dot K_{NT}K_T^{-2}.
\end{align}
These identities follow by differentiating $KA=AK=P$, $A=A^*$,
$K=K^*$, and $A(I-P)=0$. For a tangent field $Y$,
\begin{align}
\nabla^E_Z(KY)-K\nabla^E_ZY
&=P_XDK_X[Z](Y)+P_XK_XDY[Z]-K_XP_XDY[Z]\\
&=P_XDK_X[Z](Y),\\
\nabla^E_Z(AK)&=0,\\
\nabla^E_ZA&=-A(\nabla^E_ZK)A .
\end{align}
Finally $\on{A|_{T_X}}\leq a^{-1}$ and
Eq.~\eqref{eq:rigid-DK} imply the bound.
\end{proof}

\subsection{Levi--Civita connection and an explicit Hessian bound}
\begin{definition}
For tangent fields $Y,Z$, define
\begin{align}
B_X(Z,Y)&=(I-P_X)DY[Z]=DP_X[Z](Y),\\
\mathfrak D_X(Z,Y)
&=\nabla^g_ZY-\nabla^E_ZY,\\
\ii{\mathfrak Q_X(Z,Y)}W
&=\ii Z{(\nabla^E_WA)Y},\qquad W\in T_X\mathcal M_r .
\end{align}
\end{definition}

\begin{theorem}[Closed connection tensor]
\label{thm:levi-civita}
\begin{align}
2g_X(\mathfrak D_X(Z,Y),W)
&=\ii Y{(\nabla^E_ZA)W}
+\ii Z{(\nabla^E_YA)W}
-\ii Z{(\nabla^E_WA)Y},\\
\mathfrak D_X(Z,Y)
&=\tfrac12 K_X\bigl[
(\nabla^E_ZA)Y+(\nabla^E_YA)Z-\mathfrak Q_X(Z,Y)\bigr],\\
\nn{B_X(Z,Y)}&\leq\frac4a\nn Z\nn Y,\\
\nn{\mathfrak D_X(Z,Y)}
&\leq\frac{12b^2}{a^3}\nn Z\nn Y .
\label{eq:rigid-connection-tensor}
\end{align}
\end{theorem}
\begin{proof}
For $\nabla^g=\nabla^E+\mathfrak D$, metric compatibility and zero torsion give
\begin{align}
Zg(Y,W)
&=g(\nabla^E_ZY,W)+g(Y,\nabla^E_ZW)
+\ii Y{(\nabla^E_ZA)W},\\
\mathfrak D(Z,Y)&=\mathfrak D(Y,Z),\\
2g(\mathfrak D(Z,Y),W)
&=(\nabla^E_Zg)(Y,W)+(\nabla^E_Yg)(Z,W)
-(\nabla^E_Wg)(Z,Y).
\end{align}
Set $C_A=4b/a^3$. Then
\begin{align}
|\ii{\mathfrak Q(Z,Y)}W|
&\leq C_A\nn Z\nn Y\nn W,\\
\nn{\mathfrak D(Z,Y)}
&\leq\tfrac12(2b)(3C_A)\nn Z\nn Y
=\frac{12b^2}{a^3}\nn Z\nn Y .
\end{align}
\end{proof}

\subsection{The pullback cancellation and a sharp Hessian bound}
\label{app:sharp-pullback}
\begin{theorem}[Second-order pullback identity]
\label{thm:sharp-pullback}
Let $p=(B,A)$ be balanced, $L_p=D\pi_p$, and let $\xi,\zeta$
be the minimum-energy horizontal lifts of $Z,Y$.
For $G=\nabla f(X)$,
\begin{equation}
{\rm Hess}_g f(X)[Z,Y]
=\ii{L_p\xi}{\nabla^2f(X)[L_p\zeta]}
+\ii G{\xi_B\zeta_A+\zeta_B\xi_A}.
\label{eq:sharp-pullback}
\end{equation}
If $k_-,k_+$ are the smallest and largest tangent eigenvalues of
$K_X$, and
$\ell I\preceq\nabla^2f(X)\preceq L_fI$ with $\ell\geq0$, then
\begin{equation}
(\ell k_- -\|G\|_2)I
\preceq\qHess_g f(X)
\preceq(L_fk_+ +\|G\|_2)I .
\label{eq:sharp-Hessian-sector}
\end{equation}
More generally,
\begin{equation}
\|\qHess_g f(X)\|_{g\to g}
\leq k_+\|\nabla^2 f(X)\|_{F\to F}+\|G\|_2
\leq2bL_f+\|G\|_2 .
\label{eq:sharp-Hessian-bound}
\end{equation}
\end{theorem}
\begin{proof}
Proposition~\ref{prop:second-order-retraction} provides a second-order
retraction. The pullback-Hessian theorem
\citep[Propositions 5.44--5.45]{boumal2023introduction} therefore gives
\begin{align}
{\rm Hess}_g f[Z,Z]
&=\left.\frac{d^2}{dt^2}
f((B+t\xi_B)(A+t\xi_A))\right|_{t=0}\\
&=\ii{L_p\xi}{\nabla^2f[L_p\xi]}
+2\ii G{\xi_B\xi_A}.
\end{align}
Polarization proves the bilinear identity. Spectral--nuclear
duality and the arithmetic--geometric mean inequality give
\begin{align}
|2\ii G{\xi_B\xi_A}|
&\leq2\|G\|_2\|\xi_B\xi_A\|_*\\
&\leq2\|G\|_2\|\xi_B\|_F\|\xi_A\|_F
\leq\|G\|_2\|\xi\|^2,\\
k_-\|\xi\|^2
&\leq\|L_p\xi\|_F^2\leq k_+\|\xi\|^2,&
\|\xi\|^2&=\|Z\|_g^2 .
\end{align}
Apply Rayleigh--Ritz and Weyl's eigenvalue perturbation inequality
\citep[Chapter III]{bhatia1997matrix}.
The connection and second-fundamental-form terms in the product
coordinate formula are combined before taking norms. Their
cancellation removes the inverse-rank-margin factors from this
Hessian bound.
\end{proof}

\begin{proposition}[Sharpness of the residual coefficient]
The coefficient one multiplying $\|G\|_2$ in
Eq.~\eqref{eq:sharp-Hessian-bound} cannot be uniformly decreased.
\end{proposition}
\begin{proof}
Take $r=1$, $m=d=2$, $B=A^\top=\sqrt\sigma\,e_1$, and
$G=\tau e_2e_2^\top$. The unit horizontal lift
\begin{equation}
\xi_B=e_2/\sqrt2,\qquad \xi_A=\pm e_2^\top/\sqrt2
\end{equation}
satisfies $2\ii G{\xi_B\xi_A}=\pm\tau=\pm\|G\|_2$.
For a linear ambient objective, its full quotient Hessian therefore
attains both endpoints of the residual bound.
\end{proof}


\begin{theorem}[Hessian: coordinate formula and pullback bound]
\label{thm:explicit-hessian}
For $G=\nabla f(X)$,
\begin{align}
\qHess_g f(X)[Z]
&=P_X\{DK_X[Z](G)+K_X(\nabla^2f(X)[Z])\}\notag\\
&\quad+\mathfrak D_X(Z,K_XG),\\
{\rm Hess}_g f(X)[Z,Y]
&=\ii Z{\nabla^2f(X)[Y]}
+\ii G{B_X(Z,Y)-\mathfrak D_X(Z,Y)} .
\label{eq:rigid-hessian-form}
\end{align}
If $\|\nabla^2f\|_{F\to F}\leq L_f$, then
\begin{equation}
\|\qHess_g f(X)\|_{g\to g}
\leq 2bL_f+\|\nabla f(X)\|_2
=:L_g .
\label{eq:rigid-Hessian-bound}
\end{equation}
\end{theorem}
\begin{proof}
\begin{align}
\qHess_g f[Z]
&=\nabla^g_Z(K_XG)
=\nabla^E_Z(K_XG)+\mathfrak D_X(Z,K_XG),\\
{\rm Hess}_g f[Z,Y]
&=Z(Df[Y])-Df[\nabla^g_ZY]\\
&=\ii Z{\nabla^2f[Y]}+\ii G{DY[Z]-\nabla^g_ZY}\\
&=\ii Z{\nabla^2f[Y]}+\ii G{B_X(Z,Y)-\mathfrak D_X(Z,Y)},\\
|{\rm Hess}_g f[Z,Y]|
&\leq(2bL_f+\|G\|_2)\gn Z\gn Y .
\end{align}
The final estimate uses Eq.~\eqref{eq:sharp-pullback} and its residual
operator bound rather than separate norms of the connection tensors.
\end{proof}

\subsection{Higher derivatives without internal singular-value gaps}
\begin{proposition}[Resolvent differentiation]
Let $\Gamma$ enclose $[a^2,b^2]$, exclude $0$, and stay in the analytic
domain of $z^{1/2}$. Put
\begin{align}
M&=XX^\top,&R(z)&=(zI-M)^{-1},&
M_Z&=ZX^\top+XZ^\top .
\end{align}
A local smooth extension agreeing with $S_L$ on $\mathcal M_r$ satisfies
\begin{align}
S_L(X)&=\frac1{2\pi{\rm i}}\oint_\Gamma z^{1/2}R(z)\,dz,\\
DS_L[Z]&=\frac1{2\pi{\rm i}}\oint_\Gamma
z^{1/2}R(z)M_ZR(z)\,dz,\\
D^2S_L[Z,Y]
&=\frac1{2\pi{\rm i}}\oint_\Gamma z^{1/2}
\bigl[R(ZY^\top+YZ^\top)R\notag\\
&\hspace{38mm}+RM_ZRM_YR+RM_YRM_ZR\bigr]\,dz .
\label{eq:rigid-resolvent}
\end{align}
These expressions require separation from the zero spectrum, not
$\sigma_i\ne\sigma_j$ for positive singular values.
\end{proposition}
\begin{proof}
\begin{align}
D(zI-M)^{-1}[Z]&=R M_ZR,\\
D^2(zI-M)^{-1}[Z,Y]
&=R(ZY^\top+YZ^\top)R+RM_ZRM_YR+RM_YRM_ZR .
\end{align}
For $\delta_\Gamma={\rm dist}(\Gamma,{\rm spec}(M))>0$,
\begin{align}
\sup_{z\in\Gamma}\on{R(z)}&\leq\delta_\Gamma^{-1},&
\on{M_Z}&\leq2b\nn Z,\\
\nn{D^2S_L[Z,Y]}
&\leq\frac{|\Gamma|\sup_\Gamma|z^{1/2}|}{2\pi}
\left(2\delta_\Gamma^{-2}
+8b^2\delta_\Gamma^{-3}\right)\nn Z\nn Y.
\end{align}
Differentiation under the contour integral follows from uniform resolvent bounds.
\end{proof}

\begin{proposition}[Coordinate metric and curvature]
For chart coordinates $\theta^i$ and $J_i=\partial_i\Psi$,
\begin{align}
g_{ij}&=\ii{J_i}{A_XJ_j},\\
\partial_k g_{ij}
&=\ii{\partial_kJ_i}{A_XJ_j}
+\ii{J_i}{DA_X[J_k]J_j}
+\ii{J_i}{A_X\partial_kJ_j},\\
\Gamma^k_{ij}
&=\tfrac12 g^{k\ell}
(\partial_i g_{j\ell}+\partial_j g_{i\ell}-\partial_\ell g_{ij}),\\
R^k{}_{\ell ij}
&=\partial_i\Gamma^k_{j\ell}-\partial_j\Gamma^k_{i\ell}
+\Gamma^k_{im}\Gamma^m_{j\ell}
-\Gamma^k_{jm}\Gamma^m_{i\ell},\\
({\rm Hess}_g f)_{ij}
&=\partial_{ij}(f\circ\Psi)-\Gamma^k_{ij}\partial_k(f\circ\Psi).
\label{eq:rigid-coordinate-curvature}
\end{align}
Repeated indices are summed.
\end{proposition}

\subsection{Balanced gradient flow and discrete imbalance}
\begin{proposition}
For $\dot B=-GA^\top$, $\dot A=-B^\top G$,
\begin{align}
\frac d{dt}(B^\top B-AA^\top)
&=-AG^\top B-B^\top GA^\top+B^\top GA^\top+AG^\top B=0,\\
\dot X&=-BB^\top G-GA^\top A=-K_XG
\quad\text{on }\mathcal S,\\
\frac d{dt}f(X)&=-\nn{GA^\top}^2-\nn{B^\top G}^2
=-\gn{\qgrad_g f}^2.
\end{align}
For a horizontal discrete step $p^+=p-\eta v$,
\begin{align}
(B^+)^\top B^+-A^+(A^+)^\top
&=\eta^2(v_B^\top v_B-v_Av_A^\top),\\
\nn{(B^+)^\top B^+-A^+(A^+)^\top}
&\leq\eta^2\nn v^2,\\
\nn{Q_{\rm bal}-I}+\nn{Q_{\rm bal}^{-1}-I}
&\leq C_{\rm bal}\eta^2\nn v^2
\end{align}
for a smooth near-identity local balancing gauge. The last bound follows
from the inverse Sylvester derivative and the implicit function theorem.
\end{proposition}

\subsection{Differentiating balance and the actual persistent-state map}
\label{app:rigid-state}
\begin{definition}[A smooth SPD balancing gauge]
For full-rank factors, set
\begin{align}
G_B&=B^\top B,&G_A&=AA^\top,\\
T&=(G_B^{1/2}G_AG_B^{1/2})^{1/2},&
P&=G_B^{-1/2}TG_B^{-1/2},&
Q&=P^{1/2},\\
\mathcal B(B,A)&=(BQ,Q^{-1}A).
\label{eq:rigid-spd-balance}
\end{align}
All square roots in this definition are positive definite.
\end{definition}

\begin{proposition}[Exact balance differential]
For $\delta p=(\delta B,\delta A)$,
\begin{align}
PG_BP&=G_A,\qquad Q^\top G_BQ=Q^{-1}G_AQ^{-\top},\\
\delta G_B&=B^\top\delta B+\delta B^\top B,\\
\delta G_A&=\delta A A^\top+A\delta A^\top,\\
(\delta P)G_BP+PG_B(\delta P)
&=\delta G_A-P(\delta G_B)P,\\
C&=G_B^{1/2}(\delta P)G_B^{1/2},\\
CT+TC&=G_B^{1/2}\{\delta G_A-P(\delta G_B)P\}G_B^{1/2},\\
Q\,\delta Q+\delta Q\,Q&=\delta P,\\
(D\mathcal B_p[\delta p])_B&=\delta BQ+B\delta Q,\\
(D\mathcal B_p[\delta p])_A&=Q^{-1}\delta A-Q^{-1}\delta Q Q^{-1}A .
\label{eq:rigid-balance-differential}
\end{align}
Thus two positive-definite Sylvester solves determine the gauge
differential without differentiating individual eigenvectors.
\end{proposition}
\begin{proof}
\begin{align}
PG_BP
&=G_B^{-1/2}T^2G_B^{-1/2}=G_A,\\
Q^2G_BQ^2=G_A
&\Longrightarrow QG_BQ=Q^{-1}G_AQ^{-1},\\
D(PG_BP)&=(\delta P)G_BP+P(\delta G_B)P+PG_B(\delta P),\\
D(Q^2)&=Q\delta Q+\delta Q Q,&
D(Q^{-1})&=-Q^{-1}\delta Q Q^{-1}.
\end{align}
Congruence of the $\delta P$ equation by $G_B^{1/2}$ gives the
Sylvester equation in $C$. Both coefficients $T,Q$ are positive
definite, so both linear equations have unique symmetric solutions.
\end{proof}

\begin{corollary}[First-order balance is a tangent projection]
At $G_B=G_A=H$,
\begin{align}
P=Q&=I,&
H\delta Q+\delta QH&=-\qsym C_p(\delta p),\\
D\mathcal B_p[\delta p]&=\delta p+V_p(\delta Q),&
D\mathcal B_p&=\operatorname{Proj}_{T_p\mathcal S},\\
\|D\mathcal B_p\|&=1,&
\delta p\in T_p\mathcal S&\Longrightarrow\delta Q=0 .
\end{align}
\end{corollary}
\begin{proof}
Since $\delta G_A-\delta G_B=-2\qsym C_p(\delta p)$ and
$\delta P=2\delta Q$, the Sylvester equation follows.
For symmetric $\delta Q$, $V_p(\delta Q)$ is normal to
$T_p\mathcal S$. Furthermore,
\begin{equation}
\qsym C_p(\delta p+V_p(\delta Q))
=\qsym C_p(\delta p)+H\delta Q+\delta QH=0 .
\end{equation}
The tangential component is unchanged and the normal component is
removed.
\end{proof}

\begin{lemma}[Differential of the horizontal projector]
For $V=V_p$, $N=V^*V$, $V^\dagger=N^{-1}V^*$, and
$\Pi=I-VV^\dagger$,
\begin{align}
\delta V(\Omega)&=(\delta B\Omega,-\Omega\delta A),\\
D\Pi_p[\delta p]
&=-\Pi(\delta V)V^\dagger-(V^\dagger)^*(\delta V)^*\Pi .
\label{eq:rigid-projector-state-derivative}
\end{align}
If $G_B,G_A\succeq\alpha I$, then
\begin{equation}
\|D\Pi_p[\delta p]\|_{\rm op}
\leq\sqrt{2/\alpha}\,\|\delta p\| .
\end{equation}
\end{lemma}
\begin{proof}
Differentiate $\Pi=I-VN^{-1}V^*$ and use
$\delta N=(\delta V)^*V+V^*(\delta V)$:
\begin{align}
\delta\Pi
&=-(\delta V)N^{-1}V^*-VN^{-1}(\delta V)^*
+VN^{-1}(\delta N)N^{-1}V^*\\
&=-\Pi(\delta V)V^\dagger-(V^\dagger)^*(\delta V)^*\Pi,\\
\|\delta V\|_{\rm op}&\leq\|\delta p\|,&
\sigma_{\min}(V)^2&=\lambda_{\min}(N)\geq2\alpha .
\end{align}
The norm bound follows from $\|V^\dagger\|\leq(2\alpha)^{-1/2}$.
\end{proof}

\begin{proposition}[Finite uniform constants on a Gram band]
Suppose
$\alpha I\preceq G_A,G_B\preceq\gamma I$ and put
\begin{align}
h&=(\gamma/\alpha)^{1/4},&
C_\Pi&=\sqrt{2/\alpha},\\
C_Q&=\frac{\gamma^{3/2}}{2\alpha^2}
(\gamma/\alpha)^{1/4}\sqrt{1+(\gamma/\alpha)^2},&
C_{\mathcal B}&=h+\sqrt\gamma\sqrt{1+h^4}\,C_Q .
\end{align}
Then
\begin{align}
\|Q\|_2,\|Q^{-1}\|_2&\leq h,&
\|DQ_p[\delta p]\|_F&\leq C_Q\|\delta p\|,\\
\|D\mathcal B_p[\delta p]\|&\leq C_{\mathcal B}\|\delta p\|,\\
\|D(T_{Q(p)}m)[\delta p,\delta m]\|
&\leq h\|\delta m\|+h^2\|m\|C_Q\|\delta p\|,
\end{align}
where $T_Qm=(m_BQ,Q^{-1}m_A)$.
For $\|G_z\|_F\leq G_f$ and
$\|\nabla^2 f_z\|_{F\to F}\leq L_f$, the fresh factor gradient obeys
\begin{align}
g_z(p)&=(G_zA^\top,B^\top G_z),\\
\|g_z(p)\|&\leq\sqrt{2\gamma}G_f=:G_{\rm fac},&
\|Dg_z(p)\|_{\rm op}&\leq2\gamma L_f+G_f=:L_{\rm fac}.
\label{eq:rigid-state-constants}
\end{align}
\end{proposition}
\begin{proof}
\begin{align}
\alpha P^2\preceq PG_BP=G_A\preceq\gamma I
&\Longrightarrow\|P\|_2\leq\sqrt{\gamma/\alpha},\\
\gamma P^2\succeq G_A\succeq\alpha I
&\Longrightarrow\lambda_{\min}(P)\geq\sqrt{\alpha/\gamma},\\
\alpha I\preceq T\preceq\gamma I,\qquad
\|\delta P\|_F
&\leq\frac{\gamma}{2\alpha^2}
\|\delta G_A-P(\delta G_B)P\|_F,\\
\|\delta G_A-P(\delta G_B)P\|_F
&\leq2\sqrt\gamma\sqrt{1+(\gamma/\alpha)^2}\|\delta p\|,\\
\|\delta Q\|_F&\leq
\frac{\|\delta P\|_F}{2\lambda_{\min}(Q)} .
\end{align}
Apply the exact differentials to obtain $C_Q,C_{\mathcal B}$.
For $L_p=D\pi_p$,
\begin{align}
Dg_z(p)[\delta p]
&=L_p^*\nabla^2f_z[L_p\delta p]
+(G_z\delta A^\top,\delta B^\top G_z),\\
\|L_p\|_{\rm op}&\leq\sqrt{2\gamma},&
\|(G_z\delta A^\top,\delta B^\top G_z)\|
&\leq G_f\|\delta p\|.
\end{align}
These imply the factor-gradient bounds.
\end{proof}

\begin{theorem}[Jacobian of the exact balanced momentum state]
\label{thm:rigid-actual-state-jacobian}
For an exact, undamped connection, use a state $(p,m)$ with
$p\in\mathcal S$ and $\Pi_pm=m$:
\begin{align}
u&=\beta m+(1-\beta)g_z(p),&
v&=\Pi_pu,&z_p&=p-\eta v,\\
Q&=Q(z_p),&p^+&=\mathcal B(z_p),&
w&=T_Qv,&m^+&=\Pi_{p^+}w .
\label{eq:rigid-canonical-state}
\end{align}
This is the first-moment iteration expressed with the next balance and
state projection performed at the end of the current step.
Its exact differential is
\begin{align}
\delta v
&=(D\Pi_p[\delta p])u+
\Pi_p\{\beta\delta m+(1-\beta)Dg_z(p)[\delta p]\},\\
\delta z_p&=\delta p-\eta\delta v,&
\delta Q&=DQ_{z_p}[\delta z_p],\\
\delta p^+&=D\mathcal B_{z_p}[\delta z_p],\\
\delta w_B&=\delta v_BQ+v_B\delta Q,&
\delta w_A&=Q^{-1}\delta v_A-Q^{-1}\delta Q Q^{-1}v_A,\\
\delta m^+
&=(D\Pi_{p^+}[\delta p^+])w+\Pi_{p^+}\delta w .
\label{eq:rigid-actual-state-jacobian}
\end{align}
\end{theorem}
\begin{proof}
Projection is linear in its vector argument and $g_z(p)$ is horizontal:
\begin{equation}
\Pi_p\{\beta m+(1-\beta)g_z(p)\}
=\beta\Pi_pm+(1-\beta)g_z(p).
\end{equation}
At the next iteration, projecting transported $v$ before or after
forming the new first moment gives the same result. This proves the
state equivalence. The differential follows successively from the
product rule, Eq.~\eqref{eq:rigid-balance-differential}, and
Eq.~\eqref{eq:rigid-projector-state-derivative}.
\end{proof}

\begin{corollary}[An explicit state sensitivity matrix]
Assume the preceding Gram band holds at $p,z_p,p^+$ and along all
connecting state and oracle interpolation paths. Let $\|m\|\leq M$ and put
\begin{align}
U&=\beta M+(1-\beta)G_{\rm fac},\\
A_v&=C_\Pi U+(1-\beta)L_{\rm fac},&
C_m&=U(C_\Pi hC_{\mathcal B}+h^2C_Q).
\end{align}
Then the componentwise Jacobian bound is
\begin{equation}
\begin{pmatrix}\|\delta p^+\|\\\|\delta m^+\|\end{pmatrix}
\leq
\underbrace{\begin{pmatrix}
C_{\mathcal B}(1+\eta A_v)&C_{\mathcal B}\eta\beta\\
hA_v+C_m(1+\eta A_v)&h\beta+C_m\eta\beta
\end{pmatrix}}_{\mathsf M}
\begin{pmatrix}\|\delta p\|\\\|\delta m\|\end{pmatrix}.
\label{eq:rigid-state-sensitivity-matrix}
\end{equation}
For $d_\lambda=\|p-p'\|+\lambda\|m-m'\|$, $\lambda>0$,
\begin{align}
d_\lambda(U_zs,U_zs')
&\leq\Lambda_\lambda d_\lambda(s,s'),\\
\Lambda_\lambda
&=\max\{\mathsf M_{11}+\lambda\mathsf M_{21},
\mathsf M_{22}+\mathsf M_{12}/\lambda\},\\
d_\lambda(U_zs,U_{z'}s)
&\leq2G_{\rm fac}(1-\beta)
\{\eta C_{\mathcal B}+\lambda(h+\eta C_m)\}.
\label{eq:rigid-explicit-state-sensitivity}
\end{align}
\end{corollary}
\begin{proof}
The exact Jacobian gives
\begin{align}
\|\delta v\|&\leq A_v\|\delta p\|+\beta\|\delta m\|,\\
\|\delta z_p\|&\leq\|\delta p\|+\eta\|\delta v\|,\\
\|\delta p^+\|&\leq C_{\mathcal B}\|\delta z_p\|,\\
\|\delta m^+\|&\leq h\|\delta v\|+C_m\|\delta z_p\|.
\end{align}
Substitution yields $\mathsf M$. Multiplying by $(1,\lambda)$ proves
the common-example bound by integration along the prescribed path.
For a changed example at a fixed state,
$\|\delta v\|\leq(1-\beta)\|g_z-g_{z'}\|
\leq2(1-\beta)G_{\rm fac}$ and
$\|\delta z_p\|\leq\eta\|\delta v\|$.
The same last two inequalities give the data perturbation bound.
\end{proof}
For $0<\beta<1$ and $\lambda=\eta\beta/(1-\beta)$, its injection
constant divided by $2\eta$ is
\begin{equation}
C_U=G_{\rm fac}\{(1-\beta)C_{\mathcal B}
+\beta(h+\eta C_m)\}.
\end{equation}
The coarse matrix $\mathsf M$ supplies finite bounds; it does not
establish $\Lambda_\lambda=1+O(\eta)$ with a step-independent constant.
Such a sharper statement must preserve the cancellations in the exact
Jacobian. In particular, no contraction follows by deleting its state
blocks.

